% SPDX-License-Identifier: Apache-2.0
%
% Formal verification in TxHDL: the tutorial on proving what a unit
% states about itself. Built by //docs:prove. Every listing is included
% from the tree and every printout, log and figure is what the build
% made by running the example and the proofs.
\documentclass[journal,9pt]{IEEEtran}

\newcommand{\listingsize}{\fontsize{6}{7}\selectfont}
% SPDX-License-Identifier: Apache-2.0
%
% The house preamble, shared by every document in this directory. Loaded
% after \documentclass, before \title. Keeping it in one file is what
% stops two articles drifting apart in font, listing style or figure
% style.
% Computer Modern. IEEEtran selects Times (ptm) by default, so the face has
% to be chosen explicitly.
%
% Latin Modern is the Computer Modern variety used here, and the choice is
% forced rather than aesthetic. Under T1 encoding, plain `cmr` has no Type 1
% outlines unless cm-super is installed, and it is not in the pinned
% distribution, so pdflatex falls back to EC bitmaps and every glyph in the
% PDF becomes a Type 3 font. That was measured: the first build produced a
% document with 10 Type 3 fonts and no Type 1. Latin Modern is Computer
% Modern redrawn with a full T1 repertoire and real .pfb outlines, so the
% same design comes out scalable.
\usepackage[T1]{fontenc}
\usepackage[utf8]{inputenc}
\usepackage{lmodern}
% microtype is not in the pinned TeX distribution. emergencystretch alone
% keeps the two-column measure from producing overfull lines.
\setlength{\emergencystretch}{3em}

\usepackage{amsmath}
\usepackage{amssymb}
\usepackage{amsthm}
\usepackage{booktabs}
\usepackage{array}
% enumitem is not in the pinned TeX distribution; the list spacing below
% does what \setlist would have done.
\usepackage{float}
\usepackage{xcolor}
\usepackage{listings}
\usepackage{tikz}
\usetikzlibrary{arrows.meta,positioning,fit,backgrounds,calc}
\usepackage[hidelinks,bookmarks=true]{hyperref}

% Numbered, definition-style box for a rule the reader is meant to apply.
\theoremstyle{definition}
\newtheorem{rulebox}{Rule}
\newtheorem{example}{Example}

% Unnumbered asides.
\theoremstyle{remark}
\newtheorem*{tip}{Tip}
\newtheorem*{pitfall}{Pitfall}

% Tighter lists, in place of enumitem's \setlist.
\setlength{\topsep}{2pt}
\setlength{\itemsep}{1pt}
\setlength{\parsep}{0pt}

% Code in the running text. The typewriter face under T1 ligates `<<`
% and `>>` into guillemets, so a shift printed as \code{<<} came out as
% a French quotation mark (issue 571). microtype's \DisableLigatures is
% not in the pinned distribution, but the primitive it rests on is
% pdfTeX's own: \pdfnoligatures strips every ligature from the font it
% is given, and the current font inside \texttt is the typewriter face
% at the size in use, so each size is stripped the first time \code
% sets it. Code wants no ligature at all: two quotes are two quotes.
\newcommand{\code}[1]{\texttt{\ifdefined\pdfnoligatures\pdfnoligatures\font\fi #1}}
\newcommand{\term}[1]{\emph{#1}}

% Syntax highlighting. Four roles, distinguishable in colour and still
% distinguishable in greyscale, because an IEEE article gets printed: the
% keyword is the only bold role, the comment is the only italic one, and the
% two remaining roles differ in lightness as well as in hue.
\definecolor{kwblue}{RGB}{20,60,150}
\definecolor{cmtgreen}{RGB}{90,120,90}
\definecolor{strred}{RGB}{150,50,40}
\definecolor{numgray}{gray}{0.45}
\definecolor{framegray}{gray}{0.70}
\definecolor{bgshade}{gray}{0.97}
% A darker fill for figure cells. The listing background has to stay very
% light to keep code readable; a figure cell has to be clearly shaded or the
% distinction it draws is invisible in print.
\definecolor{cellshade}{gray}{0.90}

% The listing size is the document's choice, because it depends on the
% column: a two-column article sets `\listingsize` to 6pt and keeps its
% listings within 70 columns, or to 5.6pt when it includes source that
% rustfmt sets at 80, which is what fits a column's frame (issue 1061);
% a one-column document keeps the body size and includes source files
% at 80. `//docs:frame_test` checks every listing line against its frame. Lines are never broken; a line that does not fit
% overflows the frame, which is visible, rather than wrapping, which is
% not.
\providecommand{\listingsize}{\scriptsize}

% `'` is deliberately not a string delimiter: the language uses it for loop
% labels, as Rust does, so treating it as a quote would swallow the rest of
% the line.
\lstdefinestyle{txhdl}{
  basicstyle=\ttfamily\listingsize,
  keywordstyle=\bfseries\color{kwblue},
  commentstyle=\itshape\color{cmtgreen},
  stringstyle=\color{strred},
  numberstyle=\tiny\color{numgray},
  morekeywords={mod,use,pub,const,type,struct,enum,trait,impl,fn,async,let,mut,
    match,if,else,loop,while,for,in,break,continue,return,as,await,
    module,bus,channel,transaction,protocol,pipeline,flow,seq,config,
    port,stage,pipe,instance,connect,bind,tag,domain,blackbox,foreign,
    property,role,signal,handler,true,false,
    sequence,interface,modport,implements,package,import,var,wire,func,proc,
    case,default,next,fork,branch,join,state,fsm},
  morecomment=[l]{//},
  morestring=[b]",
  showstringspaces=false,
  columns=fullflexible,
  keepspaces=true,
  breaklines=false,
  % Framed, so a listing is bounded rather than floating in the column.
  frame=single,
  rulecolor=\color{framegray},
  framesep=3pt,
  backgroundcolor=\color{bgshade},
  xleftmargin=3pt,
  xrightmargin=3pt,
  aboveskip=6pt,
  belowskip=6pt,
  % Every listing is numbered and captioned, so it can be referred to by
  % number. `caption` without `float` numbers the listing and leaves it
  % where it was written, which is what a listing in running prose needs.
  captionpos=b,
  abovecaptionskip=2pt,
  belowcaptionskip=6pt,
}

% Figures drawn as text. Framed the same way, and with the highlighting
% switched off, because a box-and-arrow diagram is not code and half its
% words turning blue reads as an accident. Setting a style adds keys rather
% than clearing the ones already set, so the three roles are neutralised by
% hand rather than by leaving them out. Program output is printed in this
% style too, and a netlist's assignment can run past the frame, so a line
% that does not fit is broken and the rest indented; source never is.
\lstdefinestyle{ascii}{
  basicstyle=\ttfamily\listingsize,
  keywordstyle=\ttfamily,
  commentstyle=\ttfamily,
  stringstyle=\ttfamily,
  columns=fullflexible,
  keepspaces=true,
  breaklines=true,
  breakindent=2em,
  frame=single,
  rulecolor=\color{framegray},
  framesep=3pt,
  backgroundcolor=\color{bgshade},
  xleftmargin=3pt,
  xrightmargin=3pt,
  aboveskip=6pt,
  belowskip=6pt,
}
\lstset{style=txhdl}

% House TikZ styles. `shorten` lives on the arrow styles rather than on each
% arrow, so every arrow in the document gets the gap, including ones added
% later. TikZ clips a path at a node's border, so without it an arrowhead
% butts against the box with no gap at all. Keep `node distance` well above
% twice the shorten, or the arrow shrinks to nothing.
\tikzset{
  box/.style={draw,rounded corners=2pt,align=center,inner sep=4pt,
              font=\footnotesize},
  boxf/.style={box,fill=bgshade},
  lbl/.style={font=\scriptsize\itshape,align=center},
  fl/.style={-{Stealth[length=4pt]},shorten >=2.5pt,shorten <=2.5pt},
  fld/.style={fl,dashed},
}


% The contents list: the class gives a section's number 2.75em, which
% a Roman numeral past thirty overruns onto the title, so the box is
% wider here, and a subsection's entry starts where a title does.
\makeatletter
\def\l@section#1#2{\addpenalty{\@secpenalty}\addvspace{1.0em plus 1pt}%
    \@tempdima 5.75em \begingroup \parindent \z@ \rightskip \@pnumwidth%
    \parfillskip-\@pnumwidth {\bfseries\leavevmode #1}\hfil\hbox to\@pnumwidth{\hss #2}\par%
    \endgroup}
\def\l@subsection{\@dottedtocline{2}{5.75em}{5em}}
\def\l@subsubsection{\@dottedtocline{3}{10.75em}{5.5em}}
\makeatother


\InputIfFileExists{buildstamp.tex}{}{}
\providecommand{\buildstamp}{Draft build (outside the Bazel flow).}

\title{Formal Verification in TxHDL}

\author{Filip Filmar and Dragiša Janković%
\thanks{Both authors are independent. Filip Filmar is the
corresponding author, at f@hdlfactory.com; Dragiša Janković is
reached at linkedin.com/in/dragisa-jankovic.}%
\thanks{This tutorial demonstrates formal property verification for developers
who have completed the introductory tutorial~\cite{tutorial} and require
properties to hold across all reachable states rather than single execution
traces. Listings are extracted directly from the repository \code{HDL/txhdl}
during compilation; logs and waveforms reflect live verification output. The
reference design appears in~\cite{examples}, and the project index~\cite{cover}
catalogues all companion documents.}%
\thanks{The authors used a large language model, Claude, as an
assistant in exploring concepts and in drafting and constructing
the documents and implementations; every repository commit documents
this collaboration and includes the prompt sequence verbatim.}%
\thanks{\buildstamp}}

\markboth{Formal Verification in TxHDL}{Filmar and Janković: Formal Verification in TxHDL}

\begin{document}
\maketitle

\begin{abstract}
Hardware modules in TxHDL declare safety invariants, input environment
constraints, and reachability targets directly in source code using
\code{check!}, \code{assume!}, and \code{cover!}. Native simulation exercises
these assertions along tested execution paths, while dual-simulator
co-simulation validates lowered netlists against simulation traces. This report
demonstrates formal property verification: SymbiYosys, driven by pinned Yosys
toolchains, proves decimal counter safety across all reachable states using the
built-in PDR engine in yosys-abc without external SMT solvers. Bounded model
checking with Z3 isolates and renders cycle-accurate counterexamples for
failing checks, while cover runs synthesize witness traces to reachability
goals. The proof executes as a hermetic Bazel test in approximately one second;
deliberately mutated designs fail as expected, confirming verification rigor.
\end{abstract}

\begin{IEEEkeywords}
Formal verification, bounded model checking, Property Directed Reachability,
SymbiYosys, Yosys, Z3, TxHDL.
\end{IEEEkeywords}

\section{Three Ways of Being Right}
\label{sec:p-three}

Regression testing in the repository relies primarily on execution comparison.
Bazel runs an example, captures signal traces across all nets, lowers the module
to VHDL and Verilog, and re-executes the trace under \code{nvc} and Verilator:
the emitted netlist must match the reference trace cycle by cycle. This workflow
flags discrepancies introduced during lowering. It cannot identify errors
present in both the Rust model and the netlist, nor can it evaluate inputs
unstimulated by the testbench.

A second methodology embeds declarative invariants into module sources. Macro
\code{check!(c, "why")} asserts that condition \code{c} holds whenever
evaluated; macro \code{assume!(c, "why")} formalizes operational constraints on
module inputs; macro \code{cover!(c, "why")} specifies reachability goals.
Native simulation panics when a check or assumption fails, and records visits
to cover targets. Lowered netlists incorporate these statements as immediate
assertions within clocked blocks, subjecting them to continuous verification
during co-simulation~\cite{tutorial, examples}.

The third approach applies formal verification. Netlist assertions describe
properties of a finite state machine. A formal engine proves or refutes these
properties across all possible inputs and reachable states, eliminating
reliance on input selection. The following sections demonstrate this flow on a
single module: verifying invariants, inspecting a counterexample, and reaching a
cover target.

\section{The Unit}
\label{sec:p-unit}

Module \code{ex\_formal} implements a decimal counter that advances by an
arbitrary step input: 0, 1, or 2, wrapping past 9. The module formalizes four
specifications. First, the registered value remains within valid decimal range
(0 to 9). Second, whenever wrap-around occurs, the subsequent state is 0 or 1.
Third, input steps never equal 3, an unhandled mode declared via \code{assume!}.
Fourth, the count reaches value 9, declared as a \code{cover!} target.
Listing~\ref{lst:p-unit} shows the complete implementation.

\lstinputlisting[rangebeginprefix=//\ begin\{,rangebeginsuffix=\},rangeendprefix=//\ end\{,rangeendsuffix=\},includerangemarker=false,linerange=unit-unit,caption={From \code{lib/examples/ex\_formal.rs}: the digit and its four statements.},label={lst:p-unit}]{lib/examples/ex_formal.rs}

The reference testbench steps the counter through 16 cycles with varying step
inputs, printing count values and recording cover events. Every assertion holds
across all evaluated edges. The lowering engine emits SystemVerilog assertions
enclosed within \code{`ifdef FORMAL} blocks under the reset guard
(Listing~\ref{lst:p-out}). Figure~\ref{fig:p-run} illustrates the resulting
simulation trace.

\lstinputlisting[style=ascii,caption={What \code{ex\_formal} printed when the build ran it: the run, the cover count, and the netlist the proof reads.},label={lst:p-out}]{out_formal.txt}

\begin{figure*}[t]
\centering
\input{formal_timing.tex}
\caption{The run: the digit counting by its step and wrapping past
nine. One run, sixteen edges, and the statements held on each of
them.}
\label{fig:p-run}
\end{figure*}

\section{The Proof}
\label{sec:p-proof}

Formal verification consumes the Verilog netlist with macro \code{FORMAL}
defined. SymbiYosys orchestrates the preparation pipeline: command
\code{read -formal} ingests assertions as formal cells, \code{prep} elaborates
the hierarchical model, and internal transformations synthesize clocked
immediate assertions into temporal state-machine properties.

Unbounded verification uses \code{abc pdr} (Property Directed Reachability),
integrated natively within yosys-abc from the pinned ASIC toolchain~\cite{cover}.
The engine operates without external SMT solvers. It synthesizes an inductive
invariant: a state set that contains the initial state, remains closed under all
valid transitions permitted by assumptions, and satisfies all safety checks.
When an invariant is found, safety holds indefinitely. For this decimal counter,
PDR establishes inductive closure at depth three (Listing~\ref{lst:p-proof}).
Engine outputs correspond directly to the two verified safety checks.

\lstinputlisting[style=ascii,caption={The proof's log, as the build wrote it: both checks proved in frame three, and no trace, since nothing failed.},label={lst:p-proof}]{formal_proof.log}

\section{A Broken Unit}
\label{sec:p-broken}

Verification tools must be validated against known failures. The test suite
constructs two mutated netlists expected to fail. The first mutation removes the
input assumption, allowing step values of 3: wrapping from 9 yields 2, violating
the post-wrap assertion. The second mutation tightens the digit bound from
\code{count <= 4'h9} to \code{count <= 4'h8}. Listing~\ref{lst:p-tight} displays
the failure report generated by \code{abc pdr}.

\lstinputlisting[style=ascii,caption={The tightened check under \code{abc pdr}: a failure, with no trace, since the engine is asked for a verdict and not a path.},label={lst:p-tight}]{formal_tight_proof.log}

While automated test suites require a boolean verdict, diagnosing errors
requires an explicit trace. Bounded Model Checking (BMC) with Z3 evaluates
reachability across 20 cycles, identifying the minimal execution sequence
violating the property (Listing~\ref{lst:p-search}).

\lstinputlisting[style=ascii,caption={The bounded search on the tightened check: assumptions and assertions checked step by step, the failure in step six, and the trace written.},label={lst:p-search}]{formal_tight_search.log}

The solver outputs a standard VCD waveform. Figure~\ref{fig:p-cex} visualizes
the counterexample: successive step inputs of 2, 2, 1, 2, and 2 advance the
counter from 0 to 9 in five cycles. On the sixth edge, the mutated assertion
detects 9 and fails. The solver generated this stimulus sequence autonomously.

\begin{figure}[htbp]
\centering
\input{formal_cex_timing.tex}
\caption{The counterexample the bounded search found: the steps that
take the count to nine, where a check of \code{count <= 8} fails.}
\label{fig:p-cex}
\end{figure}

\section{The Cover Point}
\label{sec:p-cover}

A cover run determines reachability: it seeks an input sequence driving the
module from reset into a state satisfying the cover predicate. The bounded
engine searches state space and produces an execution path upon satisfaction
(Listing~\ref{lst:p-cover}). Figure~\ref{fig:p-reach} illustrates the witness
trace: four steps of 2 followed by one step of 1 reach terminal count 9 in six
cycles.

\lstinputlisting[style=ascii,caption={The cover run: the count reaching nine in step six, and the trace to it.},label={lst:p-cover}]{formal_covered.log}

\begin{figure}[htbp]
\centering
\input{formal_reach_timing.tex}
\caption{The path to the cover point: four steps of two, one of one,
and the count is nine.}
\label{fig:p-reach}
\end{figure}

\section{The Rule and the Tests}
\label{sec:p-rule}

A portable shell script and two Bazel rules in \code{tools/formal} manage the
entire formal workflow. Rule \code{formal\_test} executes SymbiYosys against
netlists emitted by \code{waveform()} targets. Tests require exact exit code
matches: 0 for passing proofs, 2 for expected assertion failures. Inductive
uncertainty (\code{UNKNOWN}, exit code 4) is treated as a test failure, ensuring
unresolved properties fail CI. Rule \code{formal\_run} captures logs and
waveforms for documentation, while \code{formal\_trace} renders TikZ timing
diagrams (Listing~\ref{lst:p-build}).

\lstinputlisting[rangebeginprefix=\#\ begin\{,rangebeginsuffix=\},rangeendprefix=\#\ end\{,rangeendsuffix=\},includerangemarker=false,linerange=formal-formal,caption={From \code{docs/BUILD.bazel}: the two broken copies, the seven tests, the four runs and the two drawings.},label={lst:p-build}]{BUILD.bazel}

Script \code{tools/formal/formal.sh} coordinates toolchain paths hermetically
(Listing~\ref{lst:p-script}). It resolves unpacked Debian trees for Yosys and
Z3, wraps auxiliary Python scripts with appropriate \code{PYTHONPATH}
environments, synthesizes \code{.sby} configuration files dynamically, and
verifies exit statuses.

\lstinputlisting[style=ascii,caption={\code{tools/formal/formal.sh}, which every proof in the tree runs through.},label={lst:p-script}]{tools/formal/formal.sh}

All dependencies remain fully hermetic. SymbiYosys is fetched from upstream
release archives by cryptographic checksum and installed into
\code{//third\_party/sby}. The Python \code{click} library is provisioned via
locked wheel rules on the hermetic Python 3.12 toolchain. Z3 is packaged as a
pinned archive in \code{//third\_party/z3}. Unbounded PDR proofs execute
entirely within yosys-abc without external solvers. The seven formal tests
execute in approximately one second each during standard \code{bazel test //...}
runs.

\section{What Is Not Here}
\label{sec:p-not}

Formal properties hold subject to stated environmental assumptions. Overly
restrictive assumptions risk proving properties vacuously; the assumption-free
mutation in the test suite verifies that input constraints are necessary and
operative.

Synthesized netlists evaluate assertions strictly outside reset assertion.
Module \code{ex\_rstcheck} demonstrates why: a counter with asynchronous reset
suppresses its output during reset and asserts that valid digits appear. The
verified netlist passes (\code{//docs:rstcheck\_prove\_test}), whereas an
un-gated variant fails immediately during reset assertion. Native simulation
adopts the same discipline: while \code{set\_reset} holds, assertions and
coverage checks remain inactive.

Ongoing formal initiatives target protocol compliance for AXI-Lite peripherals
and processor correctness for Vreteno via the formal RISC-V interface (RVFI).
Both target the identical verification infrastructure demonstrated here on
larger hardware designs.

\begin{thebibliography}{9}

\bibitem{tutorial}
F.~Filmar and D.~Janković, ``TxHDL, step by step,'' project
repository \code{HDL/txhdl}, \code{//docs:tutorial}, 2026.

\bibitem{examples}
F.~Filmar and D.~Janković, ``TxHDL by example,'' project repository
\code{HDL/txhdl}, \code{//docs:examples}, 2026.

\bibitem{cover}
F.~Filmar and D.~Janković, ``TxHDL documents,'' project repository
\code{HDL/txhdl}, \code{//docs:cover}, 2026.

\bibitem{sby}
YosysHQ, ``SymbiYosys (sby), front-end for Yosys-based formal
verification flows,'' \code{github.com/YosysHQ/sby}, release
\code{v0.52}, 2025.

\end{thebibliography}

\end{document}
